% Zombieproof Medic — editable source
% The illustrated PDF and EPUB are the designed edition (Fraunces / Literata, 6x9).
% Edit this file. Blanks are \blankline{Label}. Change the words inside the braces.
% Compile with: pdflatex Zombieproof-Medic.tex
\documentclass[10pt,openright]{book}
\usepackage[utf8]{inputenc}
\usepackage[T1]{fontenc}
\usepackage[paperwidth=6in,paperheight=9in,margin=0.72in]{geometry}
\usepackage{microtype}
\usepackage{enumitem}
\setlist[itemize]{leftmargin=1.2em,itemsep=0.2em,topsep=0.35em}
\setlength{\parskip}{0.4em}
\setlength{\parindent}{1.05em}
\newcommand{\blankline}[1]{%
  \par\medskip
  {\small\scshape #1}\par\noindent
  \rule{0.8\textwidth}{0.4pt}\par\smallskip
}
\newcommand{\fieldrule}[1]{\begin{quote}\textbf{Field rule.}\enspace #1\end{quote}}
\begin{document}
\frontmatter
\begin{titlepage}
\begin{center}
{\small\scshape A field manual\par}
\vspace{1.5em}
{\LARGE Zombieproof Medic\par}
\vspace{0.7em}
{\large How to Keep the Living Living\par}
\vspace{2em}
Volume V
\end{center}
\blankline{Author}
\blankline{Publisher}
\blankline{City}
\end{titlepage}
\chapter*{Copyright}
Copyright \textcopyright{}
\blankline{Year}
\blankline{Author}
The author, the dedication, and the publisher are blank on purpose.

All rights reserved. No part of this book may be reproduced without the prior written permission of the publisher, except for brief passages in a review or a workshop notebook.

This is a work of humour and civilian first-aid education. It is not medical advice, a prescription, or permission to perform surgery, start a drip, or choose an antibiotic. Guidelines change. The dispatcher, the packet, and a clinician outrank this file. In the UK, call 999 or 112 for a life-threatening emergency. Elsewhere, use the local number. The undead are a metaphor. The bleeding is not.
\blankline{Publisher}
\blankline{Place of publication}
\blankline{ISBN}
\chapter*{Dedication}
For
\blankline{Name}
\blankline{A line}
\blankline{Another}
\chapter*{Foreword}
Someone always asks, at the bad moment, what to do. They mean the blood. They mean the person who has gone quiet. They mean the cupboard.

This edition will not turn a reader into a surgeon because the phone is sulking. The order is the civilian one: a safer scene, help, breathing, bleeding you can see, warmth, and a note the next person can read.
\blankline{Name}
\blankline{Place and date}
\mainmatter
\chapter{The Page You Read Before the Screaming}
\begin{quote}\emph{The doorway is not a treatment area until the roof has agreed.}\end{quote}
The first casualty of the apocalypse will probably be a man called Nigel who climbs through a broken pharmacy window while holding his torch between his teeth. Nigel has watched six survival films. None of them mentioned glass.

This book is for keeping Nigel, you, and the quieter people alive through disruption: a storm, evacuation, prolonged outage, isolated journey, or the unusually persistent queue outside the garden gate.

The zombies supply the scenery. Human bodies supply the rules.

There is no real-world zombie infection, and these pages offer no fictional cure disguised as medicine.

\section*{Before you become the second patient}
Stop outside the danger long enough to notice it. Traffic, unstable walls, smoke, electricity, aggressive people, floodwater, chemicals, machinery and falling objects can defeat excellent first aid. Look for a safe approach and a safe way out. A room that contains a collapsed person may also contain the reason they collapsed. You cannot see carbon monoxide by squinting with greater commitment.

Do not enter smoke, a suspected toxic atmosphere, a confined space, or water you cannot safely negotiate. Do not touch someone still in contact with an electrical supply. Turn a supply off only if you can do that safely. A generator belongs outdoors, away from doors, windows and air intakes; a cracked garage door is not a ventilation plan. If several people develop headache, dizziness or confusion in one place, leave for fresh air and seek emergency help. Do not return to investigate. Use gloves where available, protect your eyes from splashes, cover your own cuts, and wash your hands. If gloves are absent, a clean barrier can reduce contact while you act on life-threatening bleeding.

Keep your hands away from your face. Gloves are neither a substitute for handwashing nor a magic amulet that remains clean after you have scratched your nose.

\section*{Make help a job, not a wish}
In the UK, call 999 or 112 for a life-threatening emergency. Give the location first, including an access point, grid reference or landmark where useful. Say what happened, the number of patients, immediate hazards, and whether anyone is unresponsive or having difficulty breathing. Use the phone's speaker and follow the call handler.

Elsewhere, use the local emergency number. If communications fail, try a safe alternative: another handset or network, an established radio contact, a staffed aid point, or a person sent along a known safe route. Send the location and essential facts in writing. Do not leave a dependent patient without considering who will keep their airway clear, maintain pressure on bleeding, or watch for deterioration. Someone must know who went, where, and when they are expected back.

Choose a coordinator early. Their qualifications can initially be the ability to speak in sentences and count people. Assign concrete tasks: call for help; bring the first-aid bag; guide responders; keep the route clear; write the times. Point at a person and wait for acknowledgement. "Somebody do something" is a theatrical instruction with an appalling completion rate.

\section*{Consent survives the grid}
Introduce yourself, explain what you want to do, and ask an awake person for permission. Preserve clothing, privacy, warmth and dignity as far as the emergency allows. Explain before you touch.

Include the person in decisions; illness does not make them communal property. Follow local rules and professional advice when capacity or consent is uncertain.

For an unresponsive person with a suspected life-threatening emergency, seek emergency help and provide necessary first aid within your competence. Avoid unnecessary exposure and keep personal information with the care team. Do not photograph a patient for entertainment. The collapse of civilisation is not an invitation to start a very disappointing media company.

\section*{What you can prepare while shops still exist}
Learn practical first aid and bleeding control from a reputable instructor. Add paediatric training if children are part of your life.

Practise finding and opening the kit in darkness, with one hand, and while someone asks three irrelevant questions. Use a training mannequin for CPR; never practise chest compressions on a healthy person.

Keep paper copies of care plans, allergies, prescriptions, emergency contacts and essential device instructions. Agree an accessible meeting place and evacuation route. Identify people who need electricity, refrigerated medicines, mobility assistance or communication support. Their plan belongs near the beginning, before the dramatic discussion about axes.

\fieldrule{Make the scene safer, summon help, treat immediate threats, then keep reassessing. A quiet patient can be in greater danger than a loud one.}
\section*{When no ambulance can reach you yet}
Delayed care changes logistics; it does not turn an untrained person into a surgeon. Continue useful basics: control bleeding, protect breathing, prevent exposure, monitor changes, document care and keep trying to arrange a safe handover. Reconsider the route as conditions change. Ask remote clinicians for advice if any channel works. Write their instructions and repeat them back.

There are injuries that cannot be fixed with a field bag. Honest recognition of that fact helps people make better decisions about evacuation and comfort. Do not spend the only working torch hunting for a heroic procedure while nobody is keeping the patient warm.

\begin{quote}\emph{Sources: R01, R02, R15. Emergency guidance is written for UK readers; local emergency numbers and services differ.}\end{quote}

\chapter{The Person Before the Problem}
\begin{quote}\emph{Practise on the mannequin. It has excellent patience and no rib-related opinions.}\end{quote}
A first-aid bag contains useful objects. It does not contain an extra brain. That is why you need a simple order of work before the world becomes a cupboard full of falling saucepans.

Start with safety, response, help and breathing. Deal promptly with an obvious life-threatening bleed. If other people are present, allocate the jobs at the same time: one calls, one controls bleeding, one assesses the person. A memorable acronym is useful only if it gets the right action done.

Is this person responding?

Speak loudly and gently touch an adult's shoulders. Ask them to open their eyes. Do not shake someone violently or test responsiveness by causing pain. If they do not respond, call 999 immediately, preferably on speaker. Ask a helper to fetch an AED. Assess breathing while the call connects; the current UK adult guidance places the call early. Open the airway with a gentle head tilt and chin lift and check for normal breathing for no more than ten seconds. Occasional gasps, laboured irregular breaths and panting in an unresponsive person are not reassurance. Do not waste time looking for a pulse if you are a lay first aider. If breathing is absent or abnormal, or you are unsure, start CPR and follow the call handler.

\section*{Adult CPR: useful repetition}
Place the heel of one hand on the lower half of the breastbone in the centre of the chest, the other on top. Keep your shoulders over your hands and your arms straight. Press 5-6 cm deep at 100-120 compressions per minute, allowing the chest to rise fully between compressions. Minimise pauses. If the person is on a bed, begin there rather than delaying CPR to move them. If trained and able, give 30 compressions followed by two rescue breaths, each just enough to make the chest start to rise. Otherwise give continuous chest compressions. Use an AED as soon as available: switch it on, expose the chest as required, follow its pad diagrams and prompts, and ensure nobody touches the person during analysis or a shock. Resume compressions promptly when instructed. An AED saying "no shock" does not necessarily mean the person is well. Continue until trained help takes over, the person shows signs of life and breathes normally, you cannot physically continue, or the scene becomes unsafe. Swap rescuers roughly every two minutes where possible without long pauses. Children and babies need age-appropriate methods, including a strong emphasis on rescue breaths; Chapter 15 gives an emergency summary.

\section*{Unresponsive, but breathing normally}
Call for help. In a person without suspected significant trauma, the recovery position helps protect the airway: roll them onto their side with the head supported and the mouth directed downward, then keep checking their breathing. Do not prop them up and leave. If breathing becomes abnormal, roll them onto their back and begin CPR. With suspected neck or spinal injury, minimise movement and support the head in the position found. Airway and breathing take priority: if vomit or obstruction threatens breathing, move them enough to protect the airway, preferably with helpers keeping head, neck and trunk aligned. Do not force a recovery position merely to finish a checklist. Follow emergency-service instructions.

\section*{Awake is not the same as fine}
An awake person can tell you where it hurts, but begin with what could kill them: major bleeding, severe breathing difficulty, collapse, confusion, chest pain, stroke symptoms or a rapidly worsening allergic reaction. Call emergency help. Do not give food, drink or tablets to someone who is drowsy, choking, vomiting heavily or unable to swallow safely.

Once immediate threats are addressed, ask what happened and what changed. Record allergies, usual medicines, relevant illnesses and the last food or drink. Note the mechanism: a fall from height, heavy crush or penetrating injury can matter more than an unimpressive mark on the skin. Ask the person what normal looks like for them. A disability or communication difference is not evidence of a new brain injury.

\section*{Choking in an adult}
If they can cough effectively, encourage coughing and watch closely.

If they cannot speak, breathe or cough effectively, lean them forward and give up to five sharp back blows between the shoulder blades, checking after each. If unsuccessful, give up to five abdominal thrusts: stand behind, place a fist above the navel and below the breastbone, grasp it and pull sharply inward and upward. Have someone call emergency help; if alone, call once the first cycle has failed and continue alternating blows and thrusts. Chest thrusts are used instead when abdominal thrusts are unsuitable, such as later pregnancy or when you cannot encircle the abdomen; learn these in practical training. If the person becomes unresponsive, lower them safely, call for help and start CPR. Never sweep blindly with a finger.

Anyone who needed abdominal thrusts needs medical assessment afterwards. Infants require a different method.

\fieldrule{Observe, act, observe again. The patient is allowed to change their condition halfway through your excellent plan.}
\begin{quote}\emph{Sources: R01-R04, R17-R19. CPR and choking require practical training; follow dispatcher guidance when available.}\end{quote}
\section*{The checklist that is not yours}
The pile behind this book included a clinician's airway card. It names machines, suction, monitors, drugs, and a surgical airway for the moment a trained team cannot ventilate and cannot intubate. It assumes oxygen, medicines, practice, and someone who can look after the person afterwards. It also says, in the print people skip, that a surgical opening in the neck is not for children, and not for simple unconsciousness.

You are not going to learn it here. If someone is choking and can cough, let them cough. If they cannot breathe or speak, use the back blows and thrusts you were taught, and get emergency help. If they are unresponsive and not breathing normally, the airway you open is a gentle head-tilt and chin-lift, and the treatment is CPR and an AED, not a knife. If they are breathing, and you do not suspect a major injury to the neck, the recovery position is the whole ambition.

A bag-mask in untrained hands can fill a stomach and call it ventilation. Leave the advanced drawer for the people whose course had a mannequin, and an instructor who was allowed to say stop.

\section*{What the letters were for}
The card calls itself MSMAID. It is only a way for a trained team to ask, before they do something invasive, whether they have a machine that can breathe for the person, suction, a way to watch them, a plan for the airway, a line into a vein, and the drugs. They sort each of those into minimum, better, and best. A bag and mask are the minimum machine. A ventilator with oxygen is the best. You can be useful in that room without being the person who decides. Fetch the thing they name. Write the number they call out, and the time beside it. Do not squeeze a whole adult bag into someone because the bag looks disappointed half empty. A visible chest rise is the ambition. Too many breaths is its own injury.

The same card lists the moments a trained team starts thinking about a surgical opening in the neck: a face that will not let air through, a neck that is swelling, burns with a voice going hoarse, soot in the mouth, a person who can no longer protect their own airway, a blockage that back blows did not shift. Those are reasons to get help faster. They are not reasons to find a blade. The card also says that opening is not for a child, and not for someone who is merely unconscious and still breathing. If you remember one line of it, remember that.

If you are the spare pair of hands, your jobs are light. Suction that is already set up, if they ask for it. A note. Warmth. And the sentence "I have not been taught this part." That sentence has saved more people than improvisation has.

\chapter{Too Many Patients, One Nigel}
\begin{quote}\emph{An orderly queue is useful. A colour is not a diagnosis.}\end{quote}
Triage is what happens when the number of needs exceeds the help immediately available. It is not a contest in whose injury looks most suitable for the film trailer. It is a repeatable way to find urgent threats and organise care.

The supplied START references describe professional mass-casualty systems, including separate considerations for children. Those systems were designed for trained teams and particular response arrangements. This chapter explains the principles without turning a colour chart into permission to pronounce someone dead.

\section*{Begin outside the crowd}
Check whether the incident is still happening. A collision can be followed by another collision. Smoke can drift. A structure can continue falling. Keep rescuers out of hazards they cannot control.

Tell emergency services that there are multiple casualties and provide an estimated number, safe access and visible dangers. An estimate is better than silence; say when it is an estimate.

Choose an easily identified safe area for people able to move. Ask them to gather there, and appoint someone to stay with them. Being able to walk does not rule out serious injury, internal bleeding, smoke exposure or deterioration. Do not send the "walking wounded" away with congratulations and a biscuit-shaped diagnosis.

Allocate the first available helpers to obvious immediate needs: calling for help, severe bleeding, airway and breathing support, and guiding responders. Identify people who are unresponsive, struggling to breathe, bleeding heavily, or becoming confused.

Follow dispatcher instructions and use first aid within your training. If somebody is not breathing normally, use the CPR guidance in Chapter 2; a lay colour assessment must not override it.

\section*{Make one place remember}
Give each patient a simple unique identifier and record where they are. If names are unknown, "Patient 03, blue coat" works until identity can be established. Keep the identifier consistent on notes and belongings. Do not attach anything tightly around the neck or an injured limb.

Record the time, main problem, care given and who is watching. Use a notebook even if a phone works. A dead battery should not erase the fact that someone received medicine or had a tourniquet applied.

Personal details are for care, not for a public noticeboard. Keep family members and carers together where safe, especially children and people who rely on another person for communication or support. Write down where a person is moved and who accompanies them. Do not separate a child from their identifying information because the clipboard stayed on a table.

\section*{Reassessment is the actual system}
An initial impression is a snapshot, not a lifetime membership.

Recheck people waiting for care. Notice breathing, responsiveness, bleeding, warmth and changes in pain. If a person is worsening, tell the coordinator and seek a higher level of help. A patient who becomes quiet may be tired, frightened or deteriorating; find out rather than celebrating improved manners.

Do not confuse immediate priority with personal worth. Familiarity, occupation, age alone, loudness and social status are poor substitutes for a trained clinical assessment. Explain delays honestly and reassure people that they remain under observation. A waiting area needs water, shelter, sanitation, information and somebody kind enough to repeat the same answer.

\section*{A small incident you can organise}
A shed roof collapses. One person is walking and talking, one has heavy leg bleeding, and one lies still. You stop anyone entering beneath the remaining unstable roof. A helper calls emergency services and describes three casualties. Another applies direct pressure to the bleeding. You assess the unresponsive person and follow the breathing/CPR pathway. The walking person waits somewhere safe with another helper, who watches for change and writes down names.

That is not the whole of disaster medicine. It is a useful beginning. A coloured tag cannot do any of those jobs without people, a safe space and repeated attention.

\fieldrule{Sort to organise help, then reassess. Never use this book's illustrations, fictional zombies or a single colour to decide that a real person should receive no care.}
\section*{Handing over without a speech}
Give the receiving team the number of people, immediate threats, locations, identifiers, treatments and important times. Say what is uncertain. "Found unresponsive at 14:10; breathing was abnormal; CPR began immediately" is useful. "He was doing a sort of thing and then we became concerned" is less portable.

Take care of the helpers after the incident. Check for injuries, exposure, missing people and exhaustion. Replace used supplies, secure patient records and allow a brief factual review. People may feel upset later even if they seemed composed while working. There is no requirement to earn distress by having the most dramatic job.

\begin{quote}\emph{Sources: R27, R30, R31. This is organisational guidance, not a substitute for a locally approved professional triage protocol.}\end{quote}
\section*{What the colours are doing}
Professional teams sometimes use a sort called START for adults, and JumpSTART for children. You do not need their cards to do the first aid in this book. You do need to know what a colour means if someone else is using one, so you do not "correct" a system you were not taught, and so you do not treat a tag as a verdict on a person's worth.

In their hands the adult sort is blunt, and it is meant to take a minute. People who can walk are gathered and watched. Walking is not the same as well. Then, for the people who cannot walk, the team looks at breathing, at circulation, and at whether the person can follow a simple command. A colour is a queue. It is redone. It is not a diagnosis, and it is not a eulogy.

Children have their own version because a child who looks as if they are not breathing sometimes starts after the airway is opened and a few breaths are given. That is a reason not to glance at a small still person and decide the story is over. It is not a reason to skip the call for help, or to invent a paediatric protocol out of a paragraph. Do the first aid you were taught for that age. Stay. Tell the arriving team what you saw and what you did.

Do not write "dead" on a person because a diagram had a box for it. If you are the only helper, you are not running a tag system. You are calling for help, stopping the bleeding you can see, and looking after the person in front of you. Then the next one.

\chapter{Keeping the Red Stuff In}
\begin{quote}\emph{The teddy has requested a dressing proportionate to the imaginary injury.}\end{quote}
Blood is a remarkably effective internal arrangement. Once it starts redecorating the pavement, the priority is to stop the leak, summon help and prevent the patient becoming colder. Nigel's opinion about whether the wound is "probably arterial" can wait outside with the other spectators.

Life-threatening bleeding may flow continuously, spurt, pool rapidly or soak clothing and dressings. A deep wound can be serious without a theatrical spray. Internal bleeding may have little visible blood.

Major trauma, faintness, confusion, rapid breathing, pallor or clamminess after injury should prompt urgent help. Skin colour changes can be easier to see on palms, lips or the inside of the eyelids; do not rely on one sign.

\section*{Pressure first}
Call 999 or have somebody call. Expose enough of the injury to find the bleeding, using scissors where needed. Press firmly and directly on the bleeding site with a dressing or the cleanest suitable cloth available. Use gloves or a barrier if possible, but do not postpone life-saving pressure while someone searches for the deluxe packet.

Maintain the pressure. If something is embedded, do not pull it out. Press around it and build padding to stabilise it without pushing it deeper. Do not probe for fragments or remove a clot to admire your work. Once bleeding is controlled, secure a dressing without cutting off circulation and keep checking. If bleeding continues through a dressing, maintain effective pressure at the source and follow emergency-service instructions; an ever-taller pile of loose gauze is not the same as pressure.

Deep wound packing and haemostatic dressings can help in suitable wounds when used correctly. Learn hands-on techniques and the product's instructions before relying on them. Do not pack the chest, abdomen, eye or a wound you cannot safely assess. The anatomical mystery tour is cancelled.

\section*{A tourniquet has a specific job}
A commercial tourniquet is for life-threatening bleeding from an arm or leg that direct pressure does not control, or where maintaining pressure is not feasible. Apply it promptly using the manufacturer's instructions and your training. It belongs above the wound and clear of a joint; current UK first-aid guidance specifies 5-7 cm above a visible injury. Tighten until the bleeding stops and note the time. Do not loosen it periodically, hide it under a blanket, or remove it because the patient reports pain. Make it visible to the receiving clinician. A second commercial tourniquet above the first may be needed if bleeding continues. Escalate urgently and keep seeking clinical help. Improvised narrow cords, wire and cable ties can cause injury without controlling bleeding. They are not equivalent equipment. There is no safe universal timetable in this book for changing or removing a tourniquet during delayed evacuation. That decision requires clinical assessment. Record accurately, maintain warmth and breathing, and pursue handover. The absence of an ambulance does not make a guessed removal time reliable.

\section*{Why the time is not a suggestion}
The notes behind this chapter had a timetable. It is worth understanding, so that nobody uses it as a recipe.

Trained people hope, when they can, to replace a tourniquet with proper pressure on the wound itself inside about two hours. Between two hours and about six, they only try that with the person watched, help ready, and a plan for the bleeding to come straight back. After roughly six hours of a limb with no blood flow, muscle has died. Letting that blood back into the body in a rush can carry potassium and acid with it, and the heart can stop. People call this the dead-limb rule. It is a reason not to loosen the tourniquet to see what happens. It is not a reason for you to decide a limb is finished, and it is not permission to remove, convert, or "check" one because the evening got long.

Write the time it went on, where a stranger can see it. Do not hide it under a blanket. Do not take it off because the pain is impressive. Pain is the tourniquet doing a cruel and necessary job. Get the person to someone who can make the later decision. Yours is already made. It stays on until that person says otherwise.

\section*{Shock is not a mood}
Circulatory shock means tissues are not receiving adequate blood flow. Severe bleeding is one cause. The person may be weak, clammy, thirsty, confused or increasingly drowsy. Call for help, control external bleeding, and keep them still and warm. Let them lie down if tolerated, but do not force a breathing-impaired or injured person into a position that makes things worse. Do not give food, alcohol or a large drink to a severely injured person who may deteriorate or need urgent procedures. Never give anything by mouth to someone whose swallowing is unsafe. Recheck breathing and responsiveness. The blanket goes under as well as over them: cold ground is a patient, efficient thief.

\section*{Bleeding that you cannot see}
After major impact or a fall, worsening abdominal pain, a swollen abdomen, collapse, increasing breathlessness or a deteriorating person may indicate serious internal injury. Do not assume that a small skin mark makes it minor. These injuries need urgent professional care; keep the person comfortable, monitor, avoid unnecessary movement and make evacuation a priority.

For an open chest wound, current UK lay first-aid guidance says to leave it uncovered rather than improvise an airtight seal. Call emergency help. Trained responders with suitable equipment may use a specialised vented dressing and monitor it. Do not insert anything into the wound or attempt needle decompression from a book.

\fieldrule{Bleeding controlled once is not bleeding controlled forever. Recheck after movement, during transport and whenever the person changes.}
\section*{What to write}
Record when bleeding was found, what you did, whether it stopped, and any tourniquet time. Record medicines only if actually given, with name, dose, route and time. Report uncertainty explicitly. The receiving clinician needs a reliable account, not a cleaner story.

\begin{quote}\emph{Sources: R03, R20. Advanced haemorrhage procedures in military references are outside this lay guide's scope.}\end{quote}

\chapter{Wounds, Bites and Other Open Invitations}
\begin{quote}\emph{Soap: the unglamorous recruit with the best attendance record.}\end{quote}
A wound is an accidental door. Your task is to reduce what enters, protect what is underneath, and notice when the building requires more than a handyman. The shiny sewing kit in the abandoned haberdashery does not change this arrangement.

\section*{The ordinary wound, done properly}
For a small cut or graze, stop bleeding with direct pressure, wash your hands and rinse the wound with clean running tap water or safe bottled water. Clean the surrounding skin with soap and water, keeping soap and antiseptic out of the wound itself. Pat the surroundings dry and cover with a suitable sterile dressing. Change it when wet, dirty or loose.

Remove only loose, easily washed-away debris. Do not dig for glass, metal or deeply embedded dirt. A wound that stays contaminated, is deep or gaping, affects sensation or movement, or overlies a joint needs assessment. Injuries to hands and feet can conceal tendon, joint or nerve damage. The ability to wiggle a finger is useful information, not a certificate that every structure beneath it is intact.

If safe water is scarce, make obtaining it a treatment priority. A dirty rag dipped in a puddle does not improve a wound's prospects. Keep clean supplies away from dirty tools and from hands that just handled the outer packaging. Use the cleanest feasible technique; do not call it sterile unless a genuinely sterile process is maintained.

\section*{The wound you should not close}
Do not stitch, staple or superglue a deep, dirty, puncture or bite wound yourself. Closing contaminated material inside can make an injury harder to treat. A tidy surface is not proof of a tidy interior. Do not explore a wound with a blade or cut away tissue from instructions in a general field guide. Commercial skin-closure products have limited indications and their own directions. They do not convert uncertain injuries into suitable ones. Seek clinical advice about closure, tetanus protection, antibiotics and further cleaning. If care is delayed, control bleeding, cover the wound, protect it from further contamination, monitor the person and continue arranging assessment.

Avoid pouring alcohol, household disinfectant or hydrogen peroxide into wounds. Do not use petrol, ash, soil, powdered antibiotics, essential oils or mystery salves. Spirits are an especially expensive way to make a bad decision smell festive. Wound irrigation and instrument sterilisation are different tasks; Chapter 12 explains the distinction.

\section*{The kitchen is not a pharmacy, even the nice jars}
The research pile was excited about honey, moss, maggots, and a bleach solution with a person's name on it. Some of those things exist in real medicine. None of them are a jar you already own.

Medical honey is treated so it does not bring its own spores. The jar by the kettle is food. Food in a deep wound is how a small injury acquires an infection, and in a baby it can be worse than that. Do not pack a wound with honey because a paragraph said sugar is hostile to germs.

Maggots used in hospitals are a particular sterile kind, counted, contained, and put there by people who can also take them off. The ones in the bin are not volunteers. Do not introduce wildlife to a wound and call it therapy.

Dried moss turns up in old accounts of absorbent dressings. A handful from a wall is dirt with a reputation. Leave it on the wall.

There is a measured, very dilute bleach solution, buffered with a little bicarbonate, that some clinicians use on certain open wounds. It does not keep, and it burns living tissue if it is mixed strong. The bottle under the sink is not that solution. Household bleach is not one strength, it is not sterile, and a splash too much kills the flesh you were hoping to save. Do not pour it, dab it, or "just a little" it into a wound, an eye, or a drink. If a clinician gives you a named product, use that product, at their strength, for their number of days.

And the old rule still rules. Do not sew a dirty, deep, or bite wound shut to make the afternoon look finished. Closing dirt inside is how a small bad thing becomes the kind that needs a surgeon and a long apology. Gas in the tissues is not a detail. It is a reason the wound was left open.

\section*{Human and animal bites}
Real bites can damage deep tissue and introduce infection. Wash promptly with soap and warm running water, control bleeding and cover with a clean dressing. A bite that breaks skin warrants medical advice, especially on a hand, face, foot or near a joint. Deep wounds, uncontrolled bleeding or serious tissue damage need urgent care. Tell the clinician what bit, where and when. Tetanus and, depending on the animal and location, rabies assessment may be needed. Seek urgent advice after potential rabies exposure; do not wait for symptoms. Do not capture or handle a dangerous animal to prove your point. Record identifying information from a safe distance. A fictional zombie bite is a story event. In a real field guide it cannot justify amputation, cautery, abandonment, or an invented infection timetable. If a person has bitten someone, treat the actual injury and possible blood exposure. Do not retaliate against the patient, however disappointing their table manners.

\section*{Infection is a change over time}
Watch for increasing pain, spreading redness or discolouration, warmth, swelling, pus, bad smell, red streaking, fever or the person becoming unwell. On darker skin, tenderness, swelling and warmth may be more obvious than redness. Compare with the previous check and write down the change. Do not rely on a photograph alone. Rapid worsening, confusion, severe breathlessness, mottled or blue/grey skin, extreme weakness or very little urine can accompany serious infection or sepsis. Seek emergency help. A person can have sepsis without every textbook symptom, and a normal-looking temperature does not make it impossible. Cover the wound, keep the person comfortable, assess breathing and arrange urgent transfer.

\section*{Antibiotics are a decision, not a seasoning}
Different wounds involve different organisms and risks. Allergies, pregnancy, age, kidney function and interactions also affect drug choice. The remaining packet in a cupboard may be the wrong drug, wrong strength or insufficient course. Do not share prescription antibiotics or use veterinary preparations. Take a prescribed course exactly as directed and seek advice if adverse effects or deterioration occur.

Antibiotics do not remove debris, repair a tendon or replace drainage or surgery when needed. Keeping the wound protected and finding appropriate assessment are still essential. The book can explain that decision; it cannot diagnose a wound from a chapter heading.

\fieldrule{Clean gently, cover sensibly, watch changes, and seek help early. Do not turn an uncertain wound into an impressive-looking closed uncertain wound.}
\begin{quote}\emph{Sources: R05, R06, R11, R29.}\end{quote}

\chapter{Bones Have Opinions About Movement}
\begin{quote}\emph{Support the person. Decline the chair if it requires a second rescue team.}\end{quote}
If an ankle resembles a poorly assembled deckchair, do not test whether a vigorous walk will improve its outlook. Injury care begins with support, checking for other danger and resisting the temptation to put everything back by pulling it.

\section*{Suspect a fracture without demanding proof}
Pain, swelling, tenderness, bruising, loss of function or an unusual shape can suggest a broken bone. A fracture can exist without obvious deformity, and being able to move does not rule it out. Treat the injury gently and obtain medical assessment. Do not repeatedly bend a limb to discover which noise it makes.

Support it in the position found with hands, folded clothing, cushions or a sling where appropriate. Cover an open wound with a sterile dressing; do not push exposed bone back inside. Control serious bleeding without pressing on protruding bone. Remove rings and watches early if you can do so easily before swelling increases. Do not force them over an injured area.

For a makeshift support, pad hard edges, avoid pressure directly over the injury and secure without constricting. Immobilisation should make movement less painful, not create a new problem. Check the fingers or toes before and after support: warmth, colour, sensation and ability to move gently if safe. New numbness, coldness, discolouration or increasing pain needs urgent reassessment and help. Do not disturb a haemorrhage-control tourniquet to perform this check. Sprains still deserve care A twisted joint may be a sprain, a fracture or both. Protect it from further injury and rest it initially. A cold pack wrapped in cloth may reduce pain and swelling; do not put ice directly on skin or leave it against numb skin. Elevate if comfortable. Do not force weight-bearing when pain is marked, and seek assessment for inability to use the limb, severe swelling, deformity or persistent symptoms. Do not assume that an elastic bandage must be tight enough to alter someone's ancestry. If used, it should support without impairing circulation, and needs regular checks. A support that was comfortable an hour ago can become restrictive as swelling grows.

\section*{Dislocations are not a party trick}
A joint that looks displaced requires professional care. Support it in the position found and arrange assessment. Do not yank, twist or lever it into place. Blood vessels and nerves may be close to the injured structures. A video in which somebody's shoulder improved after a tug is not a reliable map of the anatomy you are currently holding.

\section*{The spine: move for a reason}
After a significant fall, collision, crush or neck/back injury, ask the person to stay still and support their head if needed. Do not straighten a painful neck or ask them to demonstrate movement. Call emergency help. Breathing and immediate danger take priority; move only as needed to maintain an airway, provide CPR, escape a genuine hazard or follow professional instructions. A coordinated move with enough helpers is preferable to several uncoordinated people pulling different corners. Explain the plan, nominate one person to give commands, and keep the person supported as a unit where feasible. Improvised rigid collars and long stays on a hard board can cause harm; do not manufacture one merely because the television ambulance had something similar.

\section*{Crush injury is more than a bruise}
A trapped limb or torso after structural collapse needs specialist rescue and urgent medical help. Continued collapse, heavy loads and sudden movement endanger rescuers as well as the trapped person.

Do not crawl beneath an unstable load. Request technical rescue, control accessible bleeding and keep the person warm and reassured while help is organised. Prolonged crushing can cause serious internal and metabolic complications when pressure is released. This is not a reason to invent a universal "leave them trapped" rule; it is a reason to involve emergency and rescue professionals as early as possible. If immediate danger forces action, follow dispatcher guidance. A general first-aid chapter cannot provide a safe release timetable or fluid regimen for every situation.

\section*{A limb that is complaining in six dialects}
A splint, a bandage, a burn, or a crush can swell until the muscle has nowhere to go. Clinicians listen for a cluster they call the six Ps. You do not need the poetry. You need the alarm.

Pain out of all proportion, and pain that screams when the limb is only gently stretched, are the early ones. The limb may feel tight, wooden, or unlike its partner. Pins and needles, a pale or dusky colour, and a pulse you can no longer feel come later, and later is a bad time to notice. On darker skin, colour is the least reliable of the six. Trust pain, tightness, and what the person says the limb is doing.

If you put the bandage or the splint on, loosen it, unless it is a tourniquet holding back life-threatening bleeding. Do not cut the skin. Do not "release the pressure" with a blade because a folder said the operation has a name. Those operations, when they are truly needed, are done by people who can also deal with the bleeding they have just created. Keep the person warm, and make this the reason the evacuation happens now.

The crush, once more, because the notes were loud about it. The dangerous moment is often not the being stuck. It is the being freed, when a long-trapped limb lets its chemistry back into the blood. Note how long they were under the thing. Do not give a large drink to someone who may need an urgent procedure. Do not follow a recipe for calcium, insulin, or bicarbonate. Those are hospital drugs, with hospital ways of killing a person who did not need them, or of saving one who did. You cannot tell which from a chapter. Your job is the rescue service, warmth, a clear note, and not crawling under a wall that has not finished falling.

\section*{Moving a patient without making a new one}
Before transport, ask: why move, to where, by which route, with whom, and what will continue during the move? Control bleeding first where possible, protect the person from weather, secure supplies and arrange a receiving contact. Take essential medicines and records. A stretcher that collapses halfway down a lane is an additional incident, not transport.

Use suitable equipment and a trained team where available. Do not lift alone if you cannot do it safely. Clear obstacles and inspect the route in advance. Agree a stop signal. Recheck breathing, bleeding, comfort and limb circulation after every move. Keep notes with the patient rather than in a vehicle travelling in the opposite direction.

\fieldrule{Support, do not repeatedly test. Move because it improves safety or access to care, not because standing still feels insufficiently heroic.}
\begin{quote}\emph{Sources: R07, R21, R22, R29.}\end{quote}

\chapter{Burns, Heads and the Eyes You Arrived With}
\begin{quote}\emph{The kettle remains a kitchen appliance, however dramatic the circumstances.}\end{quote}
The camp has survived three weeks of imaginary zombies. It now faces a kettle, a loose paving stone and a person grinding metal without goggles. Preventive medicine continues to have a frustratingly ordinary workload.

\section*{Thermal burns: cool the injury}
Make the scene safe and stop the burning process. Cool a thermal burn under cool or lukewarm running water for twenty minutes.

Remove nearby jewellery and clothing if they are not stuck to the skin. Do not pull away stuck material. Keep the rest of the person warm, particularly children, while cooling the burn. Do not use ice, butter, oils or creams.

After cooling, loosely cover with a suitable sterile non-fluffy dressing or a layer of clean cling film laid over the area, not wrapped tightly around a limb. Do not put cling film over a face. Do not burst blisters.

Burns that are large, deep, chemical or electrical, or affect important areas such as the face, hands, genitals or major joints need urgent assessment; so do burns with breathing concerns or a person becoming unwell. If clean running water is unavailable, use available clean water for cooling while arranging help; do not wait for a more elegant tap.

Avoid chilling the whole person or immersing a small child in cold water. A burn is a local injury with whole-person consequences, particularly when extensive.

\section*{Smoke and chemicals change the job}
A person with hoarseness, soot around the mouth, facial burns or breathing difficulty after smoke exposure needs emergency help.

Move to fresh air only when safe and keep monitoring. Do not enter a dangerous atmosphere to rescue without the correct equipment and training. An ordinary dust mask cannot make toxic smoke safe. For chemical exposure, protect yourself and consult the product label and emergency services. Brush dry material away carefully before irrigation when appropriate; remove contaminated clothing without spreading the chemical. Avoid contact with the substance and do not attempt chemical neutralisation. Many chemical eye or skin exposures require prompt prolonged irrigation, but unusual water-reactive materials require specific instructions. Preserve the product name for responders. Electrical injury can cause internal harm despite small skin marks.

Isolate electricity safely before approaching, call emergency help and assess response and breathing. Do not touch a person still connected to the source, and stay clear of fallen power lines. The switch labelled "probably off" has not earned your trust.

\section*{Head injury: watch the person}
After a blow to the head, loss of consciousness, seizure, worsening drowsiness, repeated vomiting, increasing headache, weakness, speech trouble, fluid from the ear or nose, or unusual behaviour warrants urgent assessment. Some people need a lower threshold for medical review, including those taking blood-thinning medicine. Tell clinicians about those medicines. A person with a minor injury still needs observation by a responsible adult for at least twenty-four hours and clear instructions about deterioration. Avoid alcohol, sedatives and return to risky work or sport. They may rest or sleep; the issue is whether they can be roused normally and whether symptoms worsen. Do not keep someone awake by force as if exhaustion were a treatment. Minimise movement after significant trauma and keep reassessing breathing. If the person becomes unresponsive, use Chapter 2. If help is delayed, record changes and keep trying to obtain urgent advice. Do not assume that "concussion" excludes bleeding inside the skull.

\section*{Eyes do not come with spares}
For a loose speck without penetrating injury, rinse gently with clean water and avoid rubbing. Do not use tweezers on the eye. Severe pain, altered vision, a high-speed fragment, an embedded object or a visible cut requires emergency assessment. Do not remove a penetrating object or press on the eyeball. Protect the area from accidental contact without pressure, and arrange urgent care. For a chemical splash, begin rinsing promptly with plenty of clean water while seeking help. Hold the eyelids open gently if possible, avoid contaminating the other eye, and remove contact lenses if easily done without delaying irrigation. NHS guidance advises at least twenty minutes of rinsing for a chemical injury. Take the chemical container or its details safely with you.

Do not use leftover prescription eye drops, household antiseptic or somebody else's contact-lens solution as treatment. Opaque fluid in a little bottle is not qualified by its bottle size.

\fieldrule{Cool a burn, protect an eye, watch a head injury, and escalate early when breathing, vision or responsiveness changes.}
\section*{The prevention shelf}
Keep hot liquids away from edges and children. Open pans away from faces. Wear suitable eye protection for cutting, grinding and chemicals. Secure lighting, clear paths and use a helmet where the task requires one. In a resource-poor setting, preventing one serious eye injury may preserve more useful work than an entire cupboard of improvised treatments.

\begin{quote}\emph{Sources: R08-R10, R15.}\end{quote}

\chapter{The Toilet Is a Life-Support System}
\begin{quote}\emph{Clean and dirty are separate departments. They do not share a bucket.}\end{quote}
There will be an enthusiastic meeting about perimeter defence.

There must also be a meeting about where the diarrhoea goes. If you only hold the first meeting, the second will eventually hold itself, throughout the settlement.

Water, hand hygiene, safe food and sewage separation are everyday medical care. They protect people before anyone reaches for a dressing. A camp that can stop a fictional horde but cannot keep faeces out of drinking water has selected a surprisingly elaborate route to failure.

\section*{Keep drinking water a separate operation}
Use a known safe supply where possible. Store drinking water in clean food-grade containers with tight lids; use a tap or pour rather than dipping hands and cups inside. Keep containers off dirty ground and away from fuels, pesticides and cleaning products. Clearly distinguish drinking water from wash water and dirty collection buckets. A container once used for chemicals does not become food-safe because the smell has become nostalgic. If microbial contamination is the concern, bring clear water to a rolling boil for one minute, or three minutes above 6,500 feet/approximately 2,000 metres. Let it cool protected and store it in a clean covered container. Settle or filter cloudy water through a clean cloth first. Boiling does not remove fuel, toxic chemicals or radioactive contamination; obtain a different source if those are suspected. A cloth filter alone does not make water microbiologically safe.

Commercial filters and disinfection tablets have different capabilities and contact times; follow the actual product's instructions and maintain it. Do not invent a universal dose for an unknown bottle of bleach. Preserve printed directions while the internet still exists.

\section*{Give hands somewhere to wash}
Put soap and safe water near food preparation and toilet areas. Wash after using the toilet, helping someone toilet, cleaning vomit or handling dirty dressings, and before preparing food, eating or giving care. Keep nails short enough to clean. Hand sanitiser is useful when appropriate, but it does not reliably replace soap and water for dirty hands or every diarrhoeal organism. Separate clean and dirty work. Store clean dressings and medicine away from washing-up, rubbish and the toilet. Use dedicated cleaning tools for contaminated areas. Gloves come off before touching communal handles. An assistant who keeps gloves on all afternoon has built a very mobile contamination network.

\section*{If toilets stop working}
Follow local emergency sanitation advice and avoid using a backed-up system. Choose a managed temporary arrangement that keeps waste contained, away from water sources and food, and inaccessible to children and animals. A fixed universal pit distance is not safe for every soil, slope, flood risk and water table. Seek public-health advice rather than digging beside the stream because the ground is softer.

Keep waste containers stable and covered, with a safe emptying or collection plan. Provide privacy, lighting, handwashing and access for people with reduced mobility. Manage menstrual products and continence supplies with equal care. Someone must own the cleaning rota, restocking and record of problems. This is a clinical role with an unfortunately small ceremonial hat.

\section*{Vomiting and diarrhoea}
The practical priority is replacing fluid losses safely. Offer frequent small sips to an awake person who can swallow, using oral rehydration solution (ORS) when indicated. Mix each sachet with exactly the stated amount of safe water; packet sizes vary. Do not make it stronger, divide a sachet by eye or add sugar. Ordinary energy drinks are not a direct substitute. Continue breastfeeding and obtain age-appropriate advice for children. Seek urgent help for inability to keep fluids down, markedly reduced urine, increasing drowsiness, blood in stool or vomit, severe abdominal pain, or a person who looks seriously unwell. Babies, older people and those with significant illness can deteriorate faster. Give nothing by mouth to a person unable to swallow safely. Keep a record of intake, vomiting, diarrhoea and urine where illness is prolonged. Small repeated amounts may be better tolerated than a large mug. Stop and reassess if drinking causes choking or worsens breathing. Do not force fluids into a sleepy person's mouth.

\section*{Stop one case becoming everyone's case}
Use separate towels, clean contaminated surfaces with the appropriate labelled product, and avoid food preparation while ill.

NHS advice generally requires staying away from work or school until at least forty-eight hours after the last vomiting or diarrhoea episode; healthcare, food work and some infections have additional rules. In a camp, keep recovering people away from food handling and seek public-health advice about an outbreak. Repeated illness in several people needs investigation of shared food, water and hygiene, not a more colourful explanation. Record who became ill and when without publicly blaming individuals. An outbreak log is a practical tool, not a list of people banned from tomorrow's committee.

\fieldrule{Safe water in; waste contained; hands washed between the two. The glamour department has no veto.}
\begin{quote}\emph{Sources: R12-R14. Use locally issued water and sanitation instructions when available.}\end{quote}
\section*{Do not invent the drink}
Oral rehydration salts, in the packet they were sold in, are the drink this book trusts for dehydration from diarrhoea. Follow that packet, and the age on it. Do not mix a homemade salt-and-sugar drink from memory, and do not "improve" a sports drink with spoons of salt because a card in a folder sounded confident. Too much salt is its own emergency. Too little does not do the job. If the person is drowsy, vomiting everything, or cannot drink safely, this is not a recipe problem. It is a reason to get help.

The packet works for a dull chemical reason. A little sodium and a little glucose, in the right proportion, cross the gut wall together, and water follows them in. That is the whole trick. A drink that is too salty, or thick with sugar, does the opposite. It pulls water out of the person and into the bowel, and the diarrhoea gets proud of itself. Neat fruit juice, undiluted energy drinks, and a double handful of salt are not "nearly" a rehydration salt. They are how a helpful person makes dehydration worse.

The proper low-osmolarity mix, per litre of safe water, is a weighed thing: 2.6 grams of sodium chloride, 13.5 grams of glucose, 1.5 grams of potassium chloride, and 2.9 grams of trisodium citrate. You will not hit those numbers with a kitchen spoon in the dark. Buy the sachets while the chemist is open. Use the line on the packet, and safe water, not a puddle. Nothing in this chapter is permission to put fluid into a person by any route except the mouth, and only when they can swallow safely. There is an old hospital method of giving fluid slowly by rectum when the mouth cannot be used. The wrong fluid, or the right fluid too fast, injures the lining you were trying to borrow. It is not in this book. Mouth, if they can swallow. Help, if they cannot.

\chapter{Weather Does Not Care About Your Backstory}
\begin{quote}\emph{Dry socks have done more for morale than several inspirational speeches.}\end{quote}
You have escaped the horde, crossed the ridge and established that your waterproof jacket is an aspirational document. The weather has no interest in your character arc. It wants heat, water and exposed skin, and it is prepared to work overnight.

\section*{Cold: warm the person, not your theory}
Hypothermia is dangerous cooling of the body. Shivering, clumsiness, slurred speech, confusion, drowsiness or slow breathing in a cold-exposed person need urgent attention. Shivering may cease in severe hypothermia; that is not proof of recovery. Seek emergency help.

Move into shelter if safe, insulate from the ground, replace wet clothing gently with dry layers and cover the head and body. Handle a severely cold or confused person gently and minimise unnecessary exertion. Warm the trunk gradually with dry insulation and a warm environment. Do not rub limbs, use a hot bath, apply fierce direct heat or give alcohol. Offer a warm non-alcoholic drink and sugary food only if fully awake and swallowing safely. Watch breathing. A cold person who becomes unresponsive needs the response and breathing assessment in Chapter 2 and emergency help; start CPR if breathing is abnormal. Do not assume death because someone is very cold and still. Lay triage shortcuts have no business here.

\section*{Heat: confusion is an emergency}
Heat exhaustion may cause heavy sweating, headache, nausea, dizziness, weakness and thirst. Stop activity, move to shade or a cool place, loosen clothing, cool the skin and offer suitable fluids if the person is alert and can swallow. Monitor their response. Heatstroke can involve confusion, collapse, seizure or very high temperature and is an emergency; sweating may still be present. Call emergency help and start active cooling immediately: use cool water, wet cloths and fanning, removing unnecessary clothing. Place wrapped cold packs at the neck and armpits if available. Do not give a drink to a confused or unresponsive person or use fever tablets as treatment for heatstroke. A person not improving promptly with cooling needs urgent assessment. Plan work for cooler periods, make shade, schedule rests and ensure access to water. Check each other, especially when someone is concentrating on a task. The confident person operating a shovel can be the last to notice that their brain has turned into porridge.

\section*{Feet: the transport system with toenails}
Inspect feet daily during sustained walking or wet conditions. Check heels, between toes, soles and pressure points. Clean and dry gently, change socks, remove grit and deal with rubbing early. Fit footwear to the actual foot and task, allowing for swelling. A blister ignored at breakfast can reorganise the expedition by supper.

Protect an intact blister with an appropriate dressing and avoid bursting it. If it opens, wash gently, preserve the remaining skin where possible, and cover. Increasing redness, warmth, swelling, pus or spreading pain needs assessment. Diabetes, poor circulation or reduced sensation lowers the threshold for seeking help. Do not let a person with numb feet use pain as their only warning system.

Prolonged wet, cold exposure can injure feet even without freezing.

Get them dry, warm gradually and avoid more exposure; numbness, swelling, unusual colour, severe pain or loss of function needs urgent advice. Do not rub damaged tissue or apply direct intense heat.

Frozen-looking tissue also needs professional assessment and protection from further injury; avoid walking on it if alternatives exist.

\section*{Shelter is part of the prescription}
Dry bedding, adequate ventilation, safe heating and enough food support recovery. Separate the sleeping area from cooking fumes, fuels and a generator. Never bring a barbecue or combustion heater intended for outdoors inside for warmth. A candle is not a medical device; use a stable, safe light and keep fire away from oxygen, bedding and loose clothing.

Make the shelter usable for the least mobile person, not only the fastest. A toilet reached over slippery steps is an injury programme.

Place water, lights and call signals within reach. Check the person who stays quiet in bed: reduced mobility and cold can become a mutually reinforcing trap.

\fieldrule{Prepare for exposure before the patient gets cold, thirsty or confused. The best emergency blanket is often the dry layer you never allowed to get wet.}
\begin{quote}\emph{Sources: R16, R29, R33. Foot-care advice is general; seek individual advice for diabetes, circulation problems or sensory loss.}\end{quote}

\chapter{The Boring Illnesses Have Excellent Attendance}
\begin{quote}\emph{The dental department opens every morning, preferably before trouble does.}\end{quote}
An apocalypse clinic will not spend all day removing cinematic arrows. It will see toothache, diarrhoea, coughs, rashes, constipation, headaches and people whose tablets ran out. This is fortunate for the costume budget and less fortunate for anyone who only packed a machete.

\section*{Notice what is new}
Ask about symptoms, onset, medicines, allergies, exposures and what the person can no longer do. Record temperature if you have a reliable thermometer, but do not let one number override appearance or deterioration. A person who becomes confused, breathless, unable to drink, increasingly weak or difficult to wake needs urgent assessment.

Useful observations include whether the person can speak comfortably, walk as usual, pass urine, keep fluids down and remain oriented. Compare with their baseline. A familiar long-term disability is not a new symptom; a sudden change in that person's function may be important. Do not diagnose by group vote.

\section*{A tooth can close the settlement}
Regular brushing with fluoride toothpaste, gentle cleaning between teeth and limiting frequent sugary snacks remain worth doing when the calendar has become speculative. Keep toothbrushes separate and clean. Dental repair depends on skills and materials that are difficult to improvise, so prevention deserves luggage space.

A suspected dental abscess needs urgent dental care. Tooth or gum pain with facial swelling, fever or a bad taste may be infection.

Difficulty breathing or swallowing, extensive swelling, or swelling near the eye needs emergency help. Do not lance it or pull a tooth with pliers. Pain relief following the product label may help while care is arranged, but does not treat the cause.

\section*{Respiratory illness and shared air}
Give an ill person rest, appropriate fluids if safe to swallow, and space from others where feasible. Improve ventilation without chilling them. Use tissues, hand hygiene and suitable masks when indicated, and clean shared touch surfaces. In a crowded shelter, protect people at greater risk and seek public-health advice if several people become ill.

Breathing difficulty, blue/grey colour, confusion, chest pain or rapidly worsening illness requires emergency help. A pulse oximeter may add information if you know how to use it, but cold fingers, movement, nail products, device limitations and skin pigmentation can affect readings. Do not delay care because a gadget displays a comforting number. A symptom is allowed to be real before a sensor agrees.

\section*{Sepsis: the person, not just the wound}
Sepsis is a life-threatening reaction to infection and can arise from more than a skin injury. Serious concern includes confusion, fast or difficult breathing, mottled or blue/grey skin, extreme illness or very little urine. Seek emergency help and say that serious infection or sepsis is a concern. Children may show different signs, including poor feeding, unusual sleepiness or reduced responsiveness. Do not wait for every sign or a particular fever threshold. Keep checking breathing, maintain comfort and warmth without overheating, and record changes while arranging transfer. Do not force fluids or tablets if swallowing is unsafe. An old antibiotic packet is not evidence that the emergency has been handled.

\section*{Abdominal pain needs humility}
Severe or sudden abdominal pain, a rigid or very tender abdomen, blood in vomit or stool, collapse, or pain with possible pregnancy needs urgent assessment. Do not assume constipation, food poisoning or anxiety because one is familiar. Do not give random laxatives for unexplained severe pain. Observe, note vomiting and bowel/urine changes, and seek advice. For minor constipation without concerning symptoms, maintaining ordinary food, fluids and movement may help; ask a pharmacist or clinician about suitable treatment. Severe illness, dehydration and particular medicines change the decision. The useful field habit is to check context rather than treat every symptom with the same bottle.

\section*{Seizures and stroke}
During a seizure, clear hazards, cushion the head, note the time and do not restrain the person or put anything in their mouth. After movements stop, assess breathing and use the recovery position when appropriate. Call emergency help for a first seizure, one lasting over five minutes, repeated seizures without recovery, serious injury or breathing concerns. Follow a known care plan if available.

For a possible stroke, think FAST: facial droop, arm weakness, speech trouble, time to call emergency help. Sudden vision loss, severe imbalance or other abrupt neurological changes can also be important. Note when the person was last known well. Do not give food, drink or aspirin as an improvised stroke treatment; swallowing may be unsafe and the cause is not established.

\section*{Medicines for ordinary discomfort}
Keep a small selection of familiar over-the-counter medicines in original packs, suitable for the actual people present. Read the strength, age restrictions, maximum daily amount, interactions and warnings each time. Do not combine cold-and-flu products with another medicine containing the same ingredient. Paracetamol overdose can seriously damage the liver; ibuprofen is unsuitable for some people, including in certain kidney, stomach, dehydration and pregnancy situations. Ask a pharmacist when uncertain. FIELD RULE: A common symptom can still be a serious illness. A familiar tablet does not establish the diagnosis.

\section*{The medicine that has to leave them awake}
A trained team sorts pain by whether the person can still think and breathe. If they can swallow, are not in shock, and need to stay sharp, the dull answer is the ordinary pain relief already in their plan. Paracetamol. An anti-inflammatory only if it is already suitable for that person. Read the packet. Do not stack two boxes that hide the same drug. Do not hand out someone else's prescription because the name sounds muscular.

Medicines that make a person sleep, drift, or forget are a different job. They can flatten breathing, especially if the person is shocked, has been drinking, or is already on something else. Ketamine turns up in military notes because, in trained hands, with a watch on the person, it can treat severe pain without dropping the blood pressure the way some other drugs do. Those are not your hands. This book will not print the milligrams. A number without a monitor is a dare.

If you think someone has taken too much of an opioid — very sleepy or unrousable, slow or stopped breathing, pinpoint pupils, blue or grey lips — call emergency help. If you have naloxone and you were shown how to use the device you actually hold, use it, then stay. It wears off before the opioid does. The person can wake up cross and then stop breathing again. If they are breathing and you are not doing chest compressions, the recovery position is the place for them. Do not leave them "sorted" because they swore at you. Swearing is not a discharge letter.

\begin{quote}\emph{Sources: R11, R34, R38, R39, R41, R47, R48.}\end{quote}

\chapter{The Looting List, With Less Looting}
\begin{quote}\emph{Useful supplies are welcome. The crystal has been referred to the decorations committee.}\end{quote}
Nigel returns from the supply run with a surgical stapler, four crystal salt lamps and no soap. He explains that the lamps were on a special offer. The inventory meeting is brief.

In fiction, the shops are abandoned. In real emergencies, use lawful supplies, organised distribution, mutual aid and help from pharmacies or health services. Do not remove stock from an operating clinic or another person's essential care. Do not enter unstable buildings, breach locked medicine storage or handle abandoned sharps. Supplies that cost a rescuer their health are a poor bargain.

\section*{Buy usefulness before drama}
Start with the people, the hazards, the available skills and the likely delay to care. A family kit, a walking kit and a shelter's shared reserve solve different problems. The following is a planning inventory rather than a universal minimum; quantities depend on the group, replenishment and training.

Need Useful items What to check Protect the helper Nitrile gloves in useful sizes; splash protection Fit, integrity, latex sensitivity, replacement supply Stop and cover bleeding Sterile dressings, gauze, tape, trauma shears Intact packs; dressings accessible with one hand Serious limb bleeding Reputable commercial tourniquet; haemostatic gauze where trained Genuine supply, instructions, hands-on training Clean and protect Soap, safe water containers, non-adherent dressings Clean/dirty separation; water supply Support and shelter Triangular bandages, padded supports, blankets Fit, warmth, mobility, dry storage Observe and record Thermometer, watch, head torch, notebook, pencils Spare power; legible units; reliable measurements Keep people well ORS sachets, suitable personal medicines, hygiene and continence products Exact users, labels, expiry, storage needs Do not prioritise injectable drugs or advanced airway equipment for an untrained group. Equipment that requires monitoring, sterile technique or specialist rescue capability is not made useful by being small and expensive. Conversely, ordinary spectacles, hearing-aid batteries, a spacer for a prescribed inhaler and a working walking aid may be indispensable to a specific person.

\section*{Inspect the whole item}
Read the product identity, strength, instructions, lot number, expiry and storage requirements. Check the seal, wrapping, contamination and signs of water, heat, freezing or pests. Do not use loose unidentified pills or decanted fluids. A bottle labelled "probably saline" is a bottle of unresolved chemistry.

Sterile packaging must be intact and the product within its labelled conditions. A dry-looking outer carton does not prove the inner sterile barrier survived. Flood-contaminated or damaged supplies should be isolated for appropriate disposal or professional assessment. Do not promote them back to the clean shelf because the cupboard is nearly empty.

For reusable equipment, retain the manufacturer's cleaning and maintenance instructions. Check batteries, seals and consumables. A functioning thermometer is more valuable than a beautiful monitor with no compatible probe. Power budgets and charging arrangements belong in the kit plan, with safe charging away from medicines and bedding.

\section*{Pack by job, not by colour aesthetic}
Keep immediate bleeding supplies near the top, in clearly marked accessible pouches. Separate routine wound care, personal medicines, hygiene and recording equipment. Make the system understandable to the person who did not pack it. Use large plain labels, accessible fastening and an inventory in the lid. Include a second light source.

Make personal essential medicines travel with their user, while keeping a documented shared reserve where lawful and appropriate.

Do not put every critical item in one bag, one vehicle or one damp cupboard. Splitting a reserve reduces single-point loss, but each cache needs suitable conditions and a named person to check it.

\section*{A supply run has an exit condition}
Before leaving, identify what is needed, where authorised supplies are available, route hazards, transport, communication and a return time. Carry a written list, containers and a way to keep temperature-sensitive items within their own requirements. Do not wander deeper into danger because a shelf looks promising.

Return, inspect and record before distributing. Note provenance and uncertain storage history. Quarantine questionable items so they cannot drift into routine use. Replace what was actually consumed, not what the shopping memory remembers. Add essentials for sanitation, menstruation, infants, disability and ongoing illness even when the dramatic kit already looks photogenic.

\fieldrule{The best item is useful to these people, within their skills, in these conditions. It is not necessarily the item that looks most like a television hospital.}
\section*{Rotation is part of acquisition}
Use first-expiring suitable stock first and replace it before expiry when supply chains are available. Record openings where an after-opening limit applies. Check reserves on a schedule and after any flood, heat event, power failure or move. A supply store is a maintained system, not a time capsule with aspirations.

\begin{quote}\emph{Sources: R23-R25, R29. The kit table is an editorial planning framework, not a clinical protocol or a substitute for local first-aid requirements.}\end{quote}
\section*{While the chemist is still open}
Buy the dull things first. Soap. Gloves that fit. Sterile dressings. Tape. A commercial tourniquet from a supplier your first-aid course recognises, and the course itself. Oral rehydration salts in their packets. A thermometer. A notebook and a pencil. Spare spectacles. The inhaler, the adrenaline auto-injector, the insulin and the glucose that already belong to a named person, with a plan from the person who prescribed them for a week when the pharmacy cannot.

Paracetamol, for someone who can read the packet and can swallow, is the pain and fever medicine this book will name, because it does not cloud a person who still has to walk and think. The dose is the one printed for that age and that weight. It is not a dare to take more, and it is not a treatment for a dirty wound. Ibuprofen has its own packet, and its own list of people who should leave it on the shelf. If a memory and a packet disagree, the packet wins.

Do not assemble a private hospital out of antibiotics, dissociative anaesthetics, or anything that belongs in a vein. Those are for a prescriber who can see the person. An antibiotic chosen from a list in a novel is how you get the wrong drug, the wrong body, and a cupboard full of false courage. Ask a pharmacist, while you still can, what this household ought to hold for the people who live in it. Write their names on the box.

Do not raid a working clinic, a locked cabinet, or a stranger's prescription. The notes that proposed a route through a hospital basement, a broken dispensing cabinet, and a sack of ampoules were filed as a bad idea with an educated vocabulary. A field manual will not improve them. If you are a clinician inside a facility that is still a facility, you already have a protocol, a register, and a reason to sign your name. This page is not that protocol.

\section*{If someone lawful hands you a box}
Some of that essay's packing list was sensible. Keep this part, and only this part, for the day a pharmacist, a clinic, or a service puts a box in your hands.

Solids travel better than liquids. A packet of oral rehydration salts, plus water you already know is safe, beats a backpack of litre bags you have not been trained to give. Glass breaks. A broken vial is a cut and a ruined medicine. Anything that must stay cold — insulin, liquid adrenaline — does not share a pocket with the plasters and the optimism. A sterile wrap that gets wet or torn has stopped being sterile. Bag it so the wrap survives the walk. Put the bleeding kit where the hand goes first. Label the pouches in words a stranger can read at two in the morning. "Top pocket: bleeding" is a system. "Miscellaneous trauma" is a hobby.

If the cabinet is locked and the hospital is still a hospital, it is locked because those drugs are someone else's emergency, and because several of them can kill a person who was only in pain. The window in chapter one was the warning. It was not the method.

\chapter{A Pharmacy Is Not a Time Capsule}
\begin{quote}\emph{Cool, dry and labelled: three qualities the cupboard must possess in reality.}\end{quote}
The medicine cupboard is warm, damp, unlabelled and located above the cooker. Nigel calls it "future-proof". In fairness, it has already become a warning to the future.

Medicines and medical supplies have different storage requirements.

Tablets, insulin, inhalers, auto-injectors, dressings and test strips are not interchangeable forms of a single substance called medical. Read the actual product and keep its instructions with it.

\section*{The long-term store}
Keep products in original labelled packaging, protected from children, unauthorised use, damp, direct sunlight, pests and excessive temperatures. Store away from sinks, cooking appliances, fuel and cleaning chemicals. Raise supplies above floors that might flood. Use a clean dry enclosure appropriate to the contents, without defeating ventilation or adding moisture. A dark cellar may be damp; a buried cache may flood; an attic may overheat. "Cool" is a measured condition, not a feeling you had while carrying boxes. Use a thermometer, record significant excursions and organise the shelf so dates can be read without dismantling it. Keep manufacturer-supplied desiccants where directed and do not improvise contact between medicines and loose drying agents.

Keep three dates distinct: expiry, opening, and any permitted period after opening or preparation. Reconstituted liquids and opened products can have much shorter limits than unopened stock. Mark the opening time or date clearly. Do not pour old and new batches together to save a bottle.

\section*{The fridge stops}
Keep the door closed while finding out which products actually require refrigeration and what the interruption means for each one.

Record the time, measured temperatures if available, and exposure history. Seek product-specific advice from a pharmacist, manufacturer or clinician. An insulated carrier helps only if its contents remain within the correct range; ice can freeze and damage medicines if used carelessly. Insulin is especially important: storage depends on the particular product, container and whether it is in use. Some products permit a defined time at room temperature; that is not permission for indefinite warm storage. Avoid freezing and direct heat. Do not stop necessary insulin or substitute a different formulation on the basis of this book; obtain urgent professional advice and monitor according to the person's care plan.

An adrenaline auto-injector is not automatically a fridge item. NHS advice warns against leaving it too hot or cold, including in a fridge.

Preserve its own instructions and check expiry. Treat all temperature limits as product-specific. The blanket rule "refrigerate anything important" is excellent for damaging some important things. If power is available from a generator, run the generator safely outdoors. Temperature monitoring and a backup plan still matter; a connected extension lead does not prove a medicine refrigerator is maintaining its range.

\section*{Expiry: evidence, not folklore}
An expiry date marks the period for which the manufacturer supports quality under the labelled conditions. Poor storage can undermine a product before that date. An apparently normal tablet after expiry does not prove potency, purity or safety. The FDA advises consumers not to use expired medicines. There is real research on extended shelf life. The US Shelf-Life Extension Program tests selected products in controlled federal stockpiles and can extend particular lots after stability assessment. It does not certify household stock or medicines recovered with an unknown history. The supplied 2024 antibiotic narrative review discusses possible extensions and important uncertainties; it is research discussion, not a dosing or scavenging protocol. In a true shortage, a qualified clinician or pharmacist may have to weigh specific risks and official emergency guidance. That is a product- and patient-specific decision. This book provides no universal extra number of safe years, no "all tablets are fine" rule and no visual test for degraded medicine. Separate uncertain stock and document what is known while seeking advice.

\section*{Clean is not sterile}
Cleaning removes dirt and material. Disinfection reduces specified microorganisms under defined conditions. Sterilisation requires a validated process capable of eliminating viable microorganisms, including resistant forms. Wiping a tool, boiling it in a saucepan or holding it over a flame does not establish a validated sterile instrument for an invasive procedure. For ordinary first aid, prioritise hand hygiene, safe water and unopened single-use sterile dressings where indicated. Do not reuse needles, syringes or other single-use invasive equipment. Do not attempt home-made injectable fluids, compounded anaesthetics or surgical instruments processed by guesswork. A bottle of boiled water is useful water; it is not an intravenous product.

Follow the actual cleaning-product label for dilution, contact time and surface compatibility. Never mix bleach with acids, ammonia or other cleaning agents. Store cleaning chemicals separately and label every prepared solution. "Strong enough to smell impressive" is not a concentration.

\section*{The five-line stock record}
For each critical product, record identity and strength, quantity and lot, expiry/opening limits, storage requirements, and location/provenance. Add temperature incidents and decisions. Keep an easy route to move doubtful stock out of active use without losing its history. A person who notices a problem should be thanked, not encouraged to put it back quietly.

\fieldrule{You cannot store medical usefulness forever.}
You can preserve identity, packaging, appropriate conditions, records and a realistic replenishment plan.

\begin{quote}\emph{Sources: R23-R26, R29, R32. No individual medicine's stability has been tested for this edition.}\end{quote}
\section*{What not to bury}
Cool, dry, dark, original bottle, off the floor, away from the cooker. That is the whole architecture this edition trusts. A buried drum of mixed tablets is not a pharmacy. It is a damp guess with a lid.

Insulin, liquid adrenaline, and anything already mixed into a syrup do not enjoy a decade underground, and this book will not give them one. Solids generally sulk more slowly than liquids. That is not the same sentence as "the date on the tub is a suggestion." Laboratories sometimes extend a dated lot under measured conditions. A biscuit tin behind the boiler is not one of those laboratories. If a pack is expired, ask a pharmacist before you decide it is still a medicine. If you cannot ask, do not promote it back to duty because the cupboard looks heroic.

If the manufacturer put a little drying packet in the bottle, leave it there. Do not tip loose silica gel onto capsules to be thorough. Too dry can crack a shell, just as too damp can cake a powder. Do not store rehydration salts in the same open tin as your tablets. They drink moisture for a living.

Keep three dates, not one: the expiry on the pack, the day you opened it, and the day someone qualified told you a particular product was finished after a hot week. Use the oldest honest pack first. Throw out anything crumbly, smelly, discoloured, leaking, or living in a bottle whose name has gone on holiday. An unlabelled pill is not a surprise. It is rubbish.

When the fridge fails, shut the door, write the time, and ask about each product on its own label. Do not invent a universal extra year. Do not freeze a medicine to "save" it. Ice is not a personality. It is a way to ruin insulin and a few other things that were having a perfectly ordinary day.

\section*{The decade people want from a hole}
The same notes offered a burial plan. A pipe, a foil bag, a drying packet, an oxygen absorber, and a promise of ten to forty years. Cool, dark, and steady is a real advantage. A hole is not automatically those things. Holes flood. Frost heaves. Labels rot off. A cache you cannot find in the dark is a cache you do not have.

At home, keep the factory bottle, seal intact, its own little drying packet still inside, off the floor, out of the bathroom steam, out of the sun, and nowhere near the cooker. Write the lot and the expiry where you can read them without unpacking the street. Use the oldest pack that is still in date. That is the whole of "first expired, first out," and it does not include expired.

Do not tip an oxygen absorber onto loose tablets. Those packets warm and sweat while they work. If a pharmacist has told you an outer bag is appropriate, the absorber stays in its sleeve, away from the pills. Do not chase the last of the moisture until gelatin capsules crack. Bone-dry is another way to ruin a medicine. Do not leave rehydration salts loose in the tin with everything else. They drink damp out of the air and turn into a brick.

A government programme has, for particular lots in particular stores, extended a date after the lot was tested. That is not a blessing you can photocopy onto a biscuit tin. And the old party story about doxycycline is wrong in one direction and useless in the other. Modern doxycycline does not rot into the kidney poison people still repeat. Left too warm, or too long, it mostly just becomes less of an antibiotic. Less of an antibiotic is not a reason to take it for a redness you have named yourself. Insulin, liquid adrenaline, eye drops, and anything you have already mixed with water do not get a decade in a pipe. They get the conditions on their label, or a pharmacist, or the bin.

\chapter{The Prescription That Must Outlive Tuesday}
\begin{quote}\emph{The treatment plan belongs to the person whose name is on it.}\end{quote}
The settlement has enough bandages to mummify a modest orchestra. It has forgotten the person whose essential medicine runs out tomorrow. A medical plan that only anticipates injuries is a plan for some entirely different group of humans.

\section*{Start with the person's existing plan}
Keep a current list of medicine names, strengths, formulations, schedules, allergies and important diagnoses, plus the prescriber's and pharmacy's details. Include devices, consumables, rescue medicines and assistance the person normally needs. Ask the person what would go wrong first if supplies or electricity stopped. They often know the answer more accurately than a general checklist.

Arrange lawful continuity before a crisis: discuss travel or emergency supplies, refill timing and alternatives with a pharmacist or prescriber. Do not hoard beyond an agreed plan or quietly share prescriptions. Abruptly stopping some medicines is dangerous, and formulations with similar names are not necessarily interchangeable. Emergency replacement arrangements vary by medicine and jurisdiction. Store paper copies and an accessible digital copy, with the person's consent. Keep a travelling set with them and a secure backup. A record locked in a powered-down office is an archive, not an emergency plan.

\section*{Diabetes: sugar is not the whole story}
People using insulin may need insulin, needles or other delivery equipment, glucose monitoring, appropriate strips or sensors, fast-acting carbohydrate, and their own sick-day and hypoglycaemia instructions. Pumps and sensors add charging and backup requirements. Temperature-sensitive supplies need their own storage plan. Do not improvise insulin switching, dosing or interruption. For a person with diabetes who is awake, can swallow safely and has a low blood glucose or their usual hypo symptoms, follow their care plan. Current first-aid guidance uses 15-20 g of glucose or dextrose; recheck after ten to fifteen minutes and repeat if still low, then give a longer-acting carbohydrate when recovered, as advised in the plan. If they are unresponsive, having a seizure or cannot swallow, give nothing by mouth. Call emergency help, assess breathing and use the recovery position if appropriate. Use prescribed glucagon only if available and someone knows how, following its instructions.

Vomiting, abdominal pain, deep rapid breathing, drowsiness or very high glucose may indicate a different diabetic emergency; seek urgent help rather than assuming more sugar is the answer.

\section*{Asthma and breathing plans}
Keep the person's own reliever, spacer where appropriate, action plan and necessary maintenance medicines together. Inhaler types and treatment plans differ; a single generic puff count is not suitable for everyone. Help the person sit comfortably upright and use their prescribed reliever according to their plan. Call emergency help for severe difficulty, inability to speak normally, exhaustion, worsening symptoms or lack of relief. Do not delay because the person is too breathless to make a dramatic speech. A wheeze after a possible allergic trigger may be anaphylaxis rather than an ordinary asthma attack. Ask about exposure and other symptoms, and act promptly. Breathing difficulty is not the time to debate the elegance of your diagnosis.

\section*{Anaphylaxis: recognise, inject, call}
Rapid swelling of the tongue or throat, breathing difficulty, wheeze, faintness or collapse after a possible trigger can be anaphylaxis. A rash may be absent. Use the person's adrenaline auto-injector promptly according to its own instructions and call 999, saying anaphylaxis is suspected. Different brands have different activation and hold instructions. Keep the person lying down; if breathing is difficult they may sit up slowly with legs extended. Do not let them stand or walk, even if they feel better. In pregnancy, lying on the left side is advised. If symptoms have not improved after five minutes, use a second auto-injector if available. Monitor breathing; start CPR if unresponsive and not breathing normally. Antihistamines do not replace adrenaline for this emergency.

\section*{Devices need an outage plan}
Oxygen concentrators, ventilators, suction equipment, feeding pumps, powered mattresses, wheelchairs and communication devices have different risks when power fails. Ask the clinical team and manufacturer for a written backup and evacuation plan. Record battery endurance under actual use and test charging safely. Do not alter oxygen flows, ventilator settings or electrical connections by guesswork.

Home oxygen increases fire risk. Keep ignition sources away and follow supplier instructions. A fuel-powered generator must stay safely outdoors; it cannot share the patient's room merely because the cable is short. If a person cannot be safely supported through the likely outage, arrange an earlier move to a suitable setting rather than waiting for the battery alarm to become the transport coordinator.

\section*{Continuity includes the unglamorous items}
Glasses, hearing-aid batteries, stoma supplies, catheters prescribed for an established user, continence products, accessible food and communication aids are part of care. Do not introduce an invasive device to solve a staffing inconvenience. Preserve existing specialist instructions and seek support when complications arise.

\fieldrule{Build the medical reserve around actual people.}
Someone's essential daily treatment outranks another person's hypothetical dramatic injury.

\begin{quote}\emph{Sources: R03, R25, R26, R35, R36, R46.}\end{quote}

\chapter{Nursing: The Bit Between Being Rescued and Being Better}
\begin{quote}\emph{The rescue is over. The care is not.}\end{quote}
The dramatic bandage has been applied. The patient has been carried through the gate. Everyone congratulates themselves and goes to sleep. Unfortunately, the patient continues to require a body for the next several days.

Prolonged care is repeated attention to breathing, warmth, comfort, skin, hydration, nutrition, medicines, hygiene and change. The field-care references devote serious space to nursing and documentation because these are treatments, not background scenery.

\section*{A place where care can happen}
Choose shelter with safe access, adequate light, ventilation, manageable temperature, handwashing and a nearby toilet or suitable alternative. Separate clean supplies from waste. Keep a call bell or another usable signal within reach, plus glasses, communication aids and essential personal items. A patient who cannot reach the water has not been supplied with water.

Arrange help before moving a person, washing them or changing bedding. Explain what will happen and preserve privacy. Keep exits clear. Avoid trailing cables, unstable chairs, open flames and loose floor coverings. The patient should not have to survive the ward furniture after surviving the roof.

\section*{Look, listen, count, write}
Record a baseline and repeat checks according to condition and clinical advice. Useful observations include responsiveness, breathing effort and rate, pulse if you know how to measure it, temperature, pain, skin changes, intake, urine and wound/dressing condition. Count breaths without asking the person to perform a special demonstration; record the number per minute and the time. A trend is more useful than an isolated confident adjective.

Write what you observed, not a diagnosis you cannot establish. "More drowsy than at 10:00; woke to voice; breathing appears harder" is useful. "Probably just tired" can erase the reason for escalation. A worsening patient needs urgent help and reassessment of immediate threats, not a more decorative chart. Do not install a universal observation interval copied from a stable patient's notes. A changing patient may need continuous attendance, while a recovering person may have a different agreed schedule.

Arrange who checks, when, and what changes require immediate action. Silence is not a completed observation.

\section*{Skin and pressure}
People unable to move are at risk of pressure damage, particularly over heels, hips, elbows and the tailbone. Check skin regularly and help them change position as safely tolerated, using an individual plan and suitable support surfaces. Look for persistent discolouration, unusual warmth, firmness, pain or broken skin; changes may appear differently on darker skin. Keep bedding clean, dry and reasonably smooth. Manage sweat, urine and stool promptly and gently. Do not massage damaged pressure areas or drag a person across bedding. Use helpers and appropriate equipment; seek advice about positioning when injuries restrict movement. A rigid two-hour rule for every patient is not a substitute for an individual pressure-relief plan.

\section*{Drinks and food without aspiration}
Offer small manageable amounts to an awake person who can swallow safely, with suitable positioning and assistance. Follow existing swallowing or feeding plans. Stop and seek advice for coughing, choking, a wet-sounding voice or difficulty managing fluids. Give nothing by mouth to someone too drowsy to protect their airway. Do not improvise feeding tubes or intravenous fluids. Provide ordinary appropriate food as tolerated when not restricted by injury or clinical advice. Include dietary needs, allergies and assistance with packaging or utensils. Monitor prolonged poor intake and weight loss where feasible. Food left outside reach is not a successful nutrition intervention.

Oral care matters even when appetite is poor. Keep the mouth clean and comfortable using the person's usual safe routine. Seek advice about painful swallowing, mouth injury or persistent inability to eat or drink. An exhausted patient should not need to deliver a formal request for a sip of water.

\section*{Dressings and medicines}
Change dressings according to the wound, product instructions and care plan, or sooner if wet, dirty, loose or compromised. Prepare supplies first, wash hands, avoid touching wound-contact surfaces and dispose of waste safely. Observe changes rather than repeatedly uncovering a clean dressing out of curiosity. Get help for infection signs or renewed bleeding.

Keep a medicine administration record with identity, medicine, dose, route, time and who gave it. Record missed or refused doses and the reason; do not quietly double the next dose. Check the current prescription and label before each administration. Separate look-alike packs and avoid interruptions during the check.

\section*{Shift changes are a treatment}
Hand over with the patient identifier, problem, allergies, usual needs, recent changes, treatment and outstanding tasks. Confirm who has responsibility next and when reassessment is due. Count critical stock and record important times. Never assume the next person heard a conversation in the doorway.

\fieldrule{A bed, a blanket and a dressing are equipment.}
Care is the continuing work that makes them useful.

\begin{quote}\emph{Sources: R27, R28, R37. The specialist military nursing references inform the principles; invasive procedures and drug regimens are not reproduced.}\end{quote}

\chapter{Humans Do Not Come in One Standard Size}
\begin{quote}\emph{An accessible table is an excellent place to begin civilisation again.}\end{quote}
The official camp plan assumes every survivor is a healthy adult who can carry twenty kilograms, climb a ladder and read small print by moonlight. This is convenient for the plan and inconvenient for almost everyone it purports to protect.

Children, pregnant people, older adults and people with disabilities or long-term conditions need preparation that starts with their actual needs. Ask them and their carers. Accessibility is easier to arrange before an emergency than while six people argue about the width of a doorway.

\section*{Children: changes can happen quickly}
Watch feeding, drinking, urine or wet nappies, interaction and breathing, not only the thermometer. A child who is unusually sleepy, difficult to wake, limp, struggling to breathe, blue/grey or markedly different from usual needs emergency assessment.

Dehydration signs include fewer wet nappies, dry mouth and reduced tears. Seek prompt medical advice for fever in a young baby; NHS guidance identifies 38°C or more below three months, and 39°C or more at three to six months, as urgent thresholds. Use medicines intended for the child's age and follow the exact product and professional directions. Different liquid strengths exist.

Use an oral syringe, not a kitchen spoon. Never estimate a fraction of an adult dose by height, and never give aspirin to a child under sixteen unless prescribed. Keep medicines secured and record who gave what, because two loving adults can accidentally give the same dose twice.

Continue usual appropriate feeding and seek help early if feeding fails. Do not dilute infant formula to stretch stock. Safe water, clean feeding equipment and an emergency feeding plan are essential; access to ready-to-feed sterile formula may be useful where appropriate. Ask a clinician or public-health team for support with infant feeding during disruption.

\section*{Baby or child unresponsive: call and act}
Make the scene safe, check response, shout for help and call 999 on speaker. Open the airway gently and check normal breathing for no more than ten seconds. If breathing is absent or abnormal, give five initial rescue breaths, each about one second with visible chest rise.

For a baby, keep the head neutral and seal over mouth and nose; for a child, use a gentle head tilt, pinch the nose and seal over the mouth. Start compressions on the lower half of the breastbone at 100-120 per minute, about one-third of chest depth: roughly 4 cm for a baby and 5 cm for a child. For infants, use two thumbs with hands encircling the chest; for children, use one or two hands as needed. A general rescuer uses 30 compressions then two breaths; specifically paediatric-trained rescuers use their 15:2 protocol. If breaths cannot be given, compressions are better than no CPR. Use an AED with paediatric mode or pads if available, following its directions; use an available AED rather than waiting for a different one. Do not let pads touch each other. Keep following the dispatcher and continue until help takes over, normal breathing and signs of life return, exhaustion prevents further action, or the scene becomes unsafe. Newborn resuscitation at birth is a separate skill and is not taught here.

If alone with no usable phone or connection, give about one minute of CPR before briefly leaving to summon help. Do not delay the call when you can make it on speaker while caring for the child. Practical paediatric first-aid training is the best preparation for this situation.

\section*{Choking in a baby}
An effectively coughing baby should be encouraged and watched, not slapped automatically. For an ineffective cough or inability to breathe, call for help. Support the baby face-down along your forearm or thigh with head lower than body, supporting the jaw without squeezing the throat, and give up to five back blows. If these fail, turn face-up with support and give up to five chest thrusts on the lower half of the breastbone using two thumbs with hands encircling the chest and fingers supporting the back. Alternate, checking after each action. Never give abdominal thrusts to an infant or sweep blindly in the mouth. If unresponsive, place on a firm surface, call emergency help and begin the baby CPR sequence. Remove only a clearly visible, easily removable object when opening the airway. Have the child assessed after significant choking treatment. Children over one year follow age-appropriate back blows and abdominal thrusts; ask the dispatcher for guidance.

\section*{Pregnancy and postpartum care}
Keep maternity records, essential medicines, contacts and an early transport plan. Bleeding, severe abdominal pain, collapse, severe headache, breathing difficulty or concern about the baby requires prompt maternity or emergency advice. Heavy bleeding, severe pain or fainting is an emergency. Do not perform internal examinations or improvised obstetric procedures from a field guide. If birth appears imminent or unexpected, call emergency services and obtain real-time guidance. Provide privacy, warmth and clean materials while help is arranged. This volume does not pretend a paragraph can replace skilled birth support or manage obstetric emergencies. Identify that need early in the community plan.

\section*{Preserve function, not just survival}
Keep mobility aids, glasses, hearing support, communication boards and familiar routines with their users. Explain care in a format the person can understand. Ask before helping or moving an aid. A wheelchair user may know exactly how they transfer safely; do not replace that knowledge with six volunteers and enthusiasm.

Older people and those with reduced sensation may show illness differently or have less reserve. A new fall, confusion, reduced intake or sudden loss of function deserves attention. Keep access to toilets, fluids, warmth and medicines dependable. Do not label a change "just age" and file the person out of view.

\fieldrule{Plan for the actual group. Equal respect often requires different equipment, timing and assistance.}
\begin{quote}\emph{Sources: R13, R14, R17-R19, R42, R43, R51, R52.}\end{quote}

\chapter{Keeping a Mind When Everything Is a Bit Much}
\begin{quote}\emph{Sometimes the useful intervention is a blanket, a quiet voice and time.}\end{quote}
After the incident, Nigel starts organising the teaspoons alphabetically. Someone else cannot stop talking. A third person wants to be alone. Stress produces a wide range of reactions; none is improved by announcing that people should simply be more resilient, preferably before dinner.

Psychological first aid is practical, humane support: make things safer, attend to immediate needs, listen without forcing a story, and connect people with useful help. It is not amateur psychotherapy, compulsory confession or a contest in optimism.

\section*{Give the next hour a shape}
Introduce yourself, speak simply and ask what the person needs now.

Offer a quieter space, water if appropriate, warmth, information and contact with someone they trust. Explain what is known and what is uncertain. Give manageable choices rather than a complete reconstruction plan for civilisation.

Do not insist that someone recount frightening events. Do not promise that everything will be fine. Do not diagnose from a reaction seen once in a crisis. Some people want conversation; others want a person nearby while they do an ordinary task. Consent and privacy apply here too. If someone is confused, first consider medical and environmental causes as well as distress: head injury, low blood sugar, infection, heat, cold, intoxication, withdrawal or carbon monoxide may be involved. New confusion needs assessment. Calling it anxiety is not an examination.

\section*{Sleep belongs in the rota}
Fatigue affects judgement, coordination and memory. Share watch and care duties so people can rest. Use a written handover rather than trusting a tired person to remember every dose, route and observation. Provide quiet periods and lighting that supports rest.

The person who says they can stay awake indefinitely has supplied a reason to check the rota, not an exemption from it.

Do not use alcohol or borrowed sedatives as treatment for distress or sleeplessness. Preserve prescribed medicines and care plans, and seek advice about missed doses, side effects or deterioration. A crisis is a poor time for abrupt unsupervised medication changes.

\section*{Panic and breathing difficulty}
Stay calm, reduce crowding and help the person settle in a comfortable position. Ask what has helped before. Encourage gentle, unforced breathing without commanding repeated deep breaths.

Never use a paper bag: the problem may not be panic, and rebreathing can be unsafe. Chest pain, fainting, unusual symptoms, severe breathlessness or a first episode deserves medical assessment. Keep explanations short. A distressed person may not process five instructions at once. One useful sentence followed by time to respond is often kinder than a lecture delivered with admirable lung capacity.

\section*{The helper is also a person}
Eat, drink, rest and hand over before exhaustion makes decisions for you. Acknowledge uncertainty and ask another person to check medicine names, records and plans. Make it acceptable to say "I need a break" or "I do not know how to do that." This is a reliability measure, not a character defect.

After difficult care, offer practical support and access to professional help. Do not require a group retelling of traumatic details. Check again later; people who seemed composed during an emergency may struggle afterwards. Confidentiality should not end when the shift ends.

\section*{When recovery is uncertain}
Some injuries and illnesses exceed what an isolated group can treat.

Continue seeking qualified advice and evacuation while providing comfort, hygiene, warmth, privacy and company. Explain honestly and gently. Respect a person's wishes and established care plans, involving professionals wherever possible.

Comfort care is active care. A supported position, clean mouth, quiet room, familiar voice and relief of avoidable discomfort matter. Use prescribed medicines only according to an appropriate plan; do not improvise sedation or end-of-life drug regimens. Clinical decisions about withholding treatment or recognising death require appropriate expertise and local processes. A fictional apocalypse does not make a colour tag a moral authority.

If someone is in immediate danger of harming themselves or another person, obtain emergency help, reduce accessible hazards when safe and stay nearby if you can do so safely. Do not physically confront an armed or violent person. The aim is protection and connection to help, not winning an argument.

\fieldrule{Give people safety, understandable information, practical help and choice. Preserve dignity when you cannot fix the situation.}
\section*{The settlement you are trying to save}
A successful group does not only retain a pulse. It retains habits that allow tired, frightened people to ask for help without humiliation.

Build predictable routines, share useful knowledge and let people contribute within their abilities. Even Nigel can check the torch batteries. In fact, he has been offered an exceptionally clear written procedure.

\begin{quote}\emph{Sources: R10, R11, R15, R35, R40, R53.}\end{quote}

\chapter{Read This When Your Brain Has Left}
\begin{quote}\emph{Leave the next person a kit they can find and a record they can read.}\end{quote}
These pages are reminders, not a replacement for the chapters or practical training. Read them while calm, keep them with the kit, and use current dispatcher or clinician instructions when available. The record sheets are intentionally blank tools for real observations; do not fill a gap with a guess. Their layouts are original for this edition, informed by R27, and are not official clinical forms.

Card A: immediate danger and adult collapse 1. Check for hazards before approaching. Do not enter smoke, toxic atmospheres, unstable structures or live electrical danger. 2. Check response. If an adult is unresponsive, call 999 or 112 immediately on speaker; send a helper for an AED. 3. Open the airway and check normal breathing for no more than ten seconds. Gasping is abnormal. 4. If not breathing normally, start CPR: centre of chest, 5-6 cm deep, 100-120 per minute, full recoil. 5. Give 30 compressions to two breaths if trained and able; otherwise continuous compressions. Use the AED and follow prompts. Nobody touches during analysis or shock. 6. If breathing normally, protect the airway and keep checking. Use recovery position when appropriate; minimise movement with major trauma unless airway or safety requires it.

Children and babies need the initial rescue breaths and age-appropriate methods in Chapter 15. Continue care until help takes over, normal breathing returns, you cannot continue or the scene becomes unsafe. [R01-R04, R51]

Card B: serious bleeding 1. Call emergency help and use firm direct pressure at the bleeding site with a dressing or clean cloth. 2. Do not remove embedded objects. Apply pressure around them without pushing them deeper. 3. For life-threatening arm or leg bleeding not controlled by pressure, use a commercial tourniquet according to training and its instructions. Tighten until bleeding stops, record time and do not release it. 4. Keep the person warm, still and monitored; recheck after movement. Give nothing by mouth if drowsy or swallowing is unsafe. 5. Pursue urgent handover. Do not stitch, cauterise or perform invasive procedures from this book.

Card C: thermal burn 1. Make the scene safe and stop the burning process.

2. Cool the burn with cool or lukewarm running water for twenty minutes. Keep the rest of the person warm. 3. Remove nearby jewellery and loose clothing; leave stuck material in place. 4. Cover loosely after cooling. No ice, butter, creams or tight wrapping; do not burst blisters. 5. Seek urgent assessment for major, deep, chemical or electrical burns, vulnerable locations, illness or breathing concerns. Chemical exposure requires specific advice.

Card D: suspected anaphylaxis 1. Rapid breathing difficulty, throat/tongue swelling, faintness or collapse after a possible trigger is an emergency.

A rash can be absent. 2. Use the person's adrenaline auto-injector promptly according to its own instructions. Call 999 and say anaphylaxis. 3. Keep lying down; if breathing is difficult, allow sitting up slowly with legs extended. Do not let them stand or walk. 4. If not improved after five minutes, use a second auto-injector if available. 5. Monitor. Start CPR if unresponsive with abnormal breathing. Antihistamines do not replace adrenaline.

Card E: diarrhoea, dehydration and infection 1. Offer frequent small sips only if awake and swallowing safely. Use suitable ORS mixed with the exact water volume on the sachet. 2. Watch intake, urine, alertness, breathing, vomiting and diarrhoea. Continue appropriate feeding and seek early advice for children or vulnerable adults. 3. Urgent help: unable to keep fluids down, markedly reduced urine, blood, severe pain or worsening illness. 4. Emergency help: confusion, difficult breathing, collapse, blue/grey or mottled skin, or difficulty waking. 5. Keep hands, safe drinking water, food and waste separated.

Several cases require an outbreak check, not blame.

Card F: before moving a patient • Why move now, and is the destination able to help?

• Is the route safe, the equipment suitable, and the team sufficient? • Is bleeding controlled as far as possible, breathing protected and the person insulated? • Who coordinates the move, watches the patient and records changes? • Are medicines, identification, care plans and records travelling with them? • Who receives the person, and what happens if the route closes?

Recheck after every move. Do not move an injured person unnecessarily, but do not leave them in immediate danger merely to avoid movement.

\section*{Patient identity and handover sheet}
Use one identifier throughout. Keep this record with the person and share only for care. Record actual observations, units, dates and times. If uncertain, write "unknown". Copy these lines into the notebook. Leave the answers blank until they are true.

- Patient identifier or name: - Date, location, and how to reach them: - Age, and any communication or access needs: - What happened, the symptoms, and when they began: - Allergies and relevant conditions: - Usual medicines and the care plan: - Major changes, and the times: - Treatment, the response, and the key times: - Last food or drink, if known: - Clinician advice, the contact, and the time: - Destination and who will receive them: - Handover: by whom, to whom, and when:

\section*{Observation sheet}
Write the real time at each check. Include units with measurements.

A deteriorating person needs action, not simply another line.

Observation frequency must suit the person and professional advice.

Time. Breathing: rate and effort. Response, pain, and skin. Action, change, and who observed.

Record temperature, pulse or other measurements below with units, device and time. Use the person's normal baseline and clinical plan where known; do not diagnose from a lone number.

- Time, measurement, units, device, and a note:

\section*{Medicine and fluid records}
Check identity, label, strength and current directions before giving medicine. Do not record a planned dose as given. Record refusals, omissions and advice sought.

- Time, medicine, strength, dose, route, who gave it, and the response: - Time, drink or food and the amount, and urine, vomit or stool:

\section*{Critical stock and cold-chain log}
Give each location a named checker. Preserve the product label and original packaging. Isolate doubtful stock; document the reason instead of silently returning it to use.

- Product, strength, lot, quantity, expiry, opened, location, and storage limits: - Date, time, temperature, how long it wandered, the advice, and who checked:

\section*{The final check before sleep}
Make the next watch explicit. Confirm who has the patient, where the medicines are, what must be checked and when to call for help.

Put the torch where the next person can find it. Restock the used dressing, not the imagined one. Write the information down.

The dead may walk. The living still need clean water, warm bedding, useful records and somebody who comes back to check. If the book has persuaded you to be that somebody, it has earned its space beside the bandages.

\section*{The Next Person Through the Door}
There is a temptation to finish a survival book by declaring that civilisation is only three meals from barbarism. This seems unfair to meals, which have generally been trying to help.

Civilisation is also somebody washing their hands before changing a dressing. Somebody making room for a wheelchair. Somebody keeping a medicine label readable, admitting that a patient is getting worse, or taking over the watch so another exhausted person can sleep.

A medical bag is useful. A group that can share it sensibly is more useful. Knowledge that can be found, understood and handed on is more useful still. Keep records so the next person does not have to guess. Preserve dignity when the answer is uncertain. Seek help early enough that a route still exists.

If the horde never arrives, all of this remains an excellent use of a Tuesday. Nigel has been placed in charge of the soap. He has a checklist. We remain cautiously optimistic.

The dead may walk. The living still need care.

\backmatter
\chapter*{Afterword}
Check them. Then quiet.
\blankline{Author}
\blankline{Place and date}
\chapter*{Acknowledgments}
The author wishes to thank
\blankline{Name}
\blankline{Name}
\blankline{Name}
and the person who said, once more, call for help before you invent a procedure.
\chapter*{About the author}
\blankline{Name}
\blankline{Born in}
\blankline{A line of biography}
\blankline{They live in}
\chapter*{Colophon}
Set, in the designed edition, in Fraunces and Literata, six by nine inches.
The plates are paintings, not demonstrations of procedures.
This file is the editable text. The author, dedication, and publisher stay blank until you fill the blankline commands.
\end{document}
